%%==========================================================================
\section{Literature Review}
\label{sec:literature}
%%==========================================================================

\subsection{Mathematical programming formulations}

Three main modelling paradigms underlie TAP research:

\paragraph{Set partitioning / column generation.}
The non-compact formulation enumerates \emph{strings} (feasible flight
sequences per aircraft) as 0--1 variables.  Because the number of strings is
exponential, column generation is required
\citep{lavoie1988new, minoux1984column, barnhart1998flight, sarac2006branch}.
Combining column generation with branch-and-bound (\emph{branch-and-price})
gives exact solutions but at significant algorithmic overhead.  Robust
extensions appear in \citet{borndorfer2010robust} and
\citet{froyland2013recoverable}.

\paragraph{Multi-commodity flow.}
\citet{levin1971scheduling} introduced the compact multi-commodity flow
formulation with $O(n_f^2 n_p)$ binary variables.  This remains the standard
benchmark for compact models and is the primary comparison point in
\citet{khaled2018compact}.

\paragraph{Time-space networks.}
\citet{hane1995fleet} proposed the time-space network for the \emph{fleet
assignment problem}.  Nodes represent airport--time events; arcs encode
ground waits, flights, and overnight stays.  Extensions to aircraft routing
appear in \citet{clarke1996maintenance}.  Preprocessing by node aggregation
and island isolation reduces model size
\citep{hane1995fleet, sherali2006airline}.  Daily and weekly maintenance
routing variants are studied in \citet{liang2011new} and
\citet{liang2012network}; the latter also integrates fleet assignment
decisions, but only approximate solutions are obtained via variable fixing.

\paragraph{Nonlinear / lifted formulations.}
\citet{haouari2012lifted} present a compact formulation with nonlinear
constraints that are linearised via the Reformulation--Linearisation Technique
(RLT).  The largest instances solved involve 344 flights, an order of magnitude
below the instances handled by the model of \citet{khaled2018compact}.

\subsection{Maintenance constraints}

Maintenance scheduling within aircraft routing has been studied since the
early work of \citet{clarke1997aircraft}.  The FAA mandates four check types
(A, B, C, D) varying in scope and frequency.  The column-generation framework
of \citet{barnhart1998flight} can embed type-A constraints through constrained
shortest-path pricing.  \citet{cordeau2001benders} apply Benders decomposition
to simultaneous routing and crew scheduling with non-linear maintenance
resources.  \citet{lacasse2010aircraft} consider maintenance under different
business-process models.  None of these works addresses the full four-level
hierarchy with explicit hierarchical counter resets, which is the subject of
\Cref{sec:hierarchy}.

\subsection{Heuristic approaches}

Constructive heuristics for TAP are comparatively under-studied relative to
exact methods.  \citet{gronkvist2005tail} provides a comprehensive treatment
including local-search improvements.  The Dijkstra-based and ACO-based
heuristics studied in the present paper are adapted from a space-time network
framework for the Locomotive Assignment Problem.

\subsection{Position of the present paper}

The compact formulation of \citet{khaled2018compact} represents the
state-of-the-art for exact, cyclic-constraint-free TAP.  The present work
builds directly upon it, providing: (i) a mathematical correction to
Constraint~(13), (ii) a hierarchical maintenance extension, and (iii) the
first systematic heuristic benchmark on the same instance family.
